\chapter{Time Integration}
\label{ch:time}

% local: velocity unknown in the first-order recast (no \vv in simbook)
\providecommand{\vv}{\mathbf{v}}

\chapterorient{Of the six analyses, the transient one is the odd member out: it
does not solve a problem, it \emph{marches} one. Where electrostatics maps a
family of independent solves and the driven analysis sweeps a band of independent
frequencies, the transient analysis threads a single field-state forward in time,
each step beginning exactly where the last one ended. This chapter is the
machinery of that march. We recast the second-order semi-discrete system as a
first-order one so off-the-shelf integrators apply; we explain why
electromagnetics almost always wants an \emph{implicit} integrator, and why each
implicit step therefore hides a linear solve; we present the integrators Palace
uses without drowning in their tableaux; and---the structural payoff---we show
that the whole march is one instance of the \code{fold\_solve} combinator of
\cref{ch:combinators}, with a single clean algebraic law that makes checkpointing
and restart fall out for free.}

\section{From a second-order system to a first-order march}
\label{sec:time-firstorder}

The transient analysis begins where \cref{ch:reductions-pde} left it. Discretizing
the time-domain wave equation \emph{in space only}---the method of lines---turned
Maxwell's time-domain system into a large \emph{second-order} system of ordinary
differential equations in the degree-of-freedom vector $\vs(t)$:
\begin{equation}
  \Mop\,\ddot\vs + \Cop\,\dot\vs + \Kop\,\vs = \vJ(t),
  \label{eq:time-secondorder}
\end{equation}
with $\Mop$ the $\varepsilon$-weighted mass, $\Cop$ the $\sigma$-weighted damping,
$\Kop$ the curl--curl stiffness, and $\vJ(t)$ the discretized forcing
$-\partial_t\vJ_{\text{src}}$ (\cref{sec:transient}). This is the equation of a
damped oscillator, with $\Mop$ the inertia, $\Kop$ the restoring stiffness, and
$\Cop$ the energy-bleeding damping---the mechanical analogy is exact. Our task in
this chapter is to produce the field history $\vs(t_0), \vs(t_1), \vs(t_2),\dots$
by advancing \cref{eq:time-secondorder} one time step at a time.

\subsection{Why first order}

Standard time integrators---the routines a numerical library ships---are written
for \emph{first-order} systems of the form
\begin{equation}
  \dot\vy = \vF(\vy, t),
  \label{eq:time-firstorder-abstract}
\end{equation}
``the rate of change of the state is a known function of the state and the time.''
A second-order system does not match this shape: it involves $\ddot\vs$, a
\emph{second} derivative. The remedy is the standard one, and it costs nothing but
bookkeeping. Carry the velocity as an extra unknown. Define
$\vv \equiv \dot\vs$ and stack the position and velocity into a \emph{doubled}
state vector $\vy = (\vs, \vv)$. Then \cref{eq:time-secondorder} becomes a pair of
first-order equations: the top one is the \emph{definition} $\dot\vs = \vv$, and
the bottom one is \cref{eq:time-secondorder} solved for $\ddot\vs = \dot\vv$:
\begin{equation}
  \frac{d}{dt}\begin{pmatrix}\vs\\\vv\end{pmatrix}
  =
  \underbrace{\begin{pmatrix}
    \vv\\
    \Mop^{-1}\bigl(\vJ(t) - \Cop\,\vv - \Kop\,\vs\bigr)
  \end{pmatrix}}_{\displaystyle \vF(\vy, t)},
  \qquad
  \vy = \begin{pmatrix}\vs\\\vv\end{pmatrix}.
  \label{eq:time-recast}
\end{equation}
This \emph{is} the first-order form \cref{eq:time-firstorder-abstract}: the doubled
state $\vy$ has rate $\vF(\vy, t)$, a known function of $\vy$ and $t$. The price is
that the state vector is now twice as long; the prize is that every integrator
ever written for $\dot\vy = \vF(\vy, t)$ now applies unchanged.
\Cref{fig:recast} draws the recast: one second-order box becomes a two-component
first-order box, with the velocity threaded between the two halves.

\begin{figure}[t]
\centering
\begin{tikzpicture}[font=\footnotesize, node distance=8mm]
  % --- LEFT: the second-order system ---
  \node[stage, minimum width=42mm, minimum height=15mm, align=center] (so) at (0,0)
    {\textbf{second order}\\[2pt]$\Mop\ddot\vs + \Cop\dot\vs + \Kop\vs = \vJ(t)$};
  \node[font=\scriptsize\itshape, text=simRust, above=1mm of so]
    {one $\ddot\vs$: library integrators don't fit};
  % arrow
  \node[right=14mm of so] (mid) {};
  \draw[flow] (so.east) -- node[above, font=\scriptsize]{$\vv \equiv \dot\vs$}
    node[below, font=\scriptsize]{recast} ($(so.east)+(2.0,0)$);
  % --- RIGHT: the first-order doubled system ---
  \node[stage, minimum width=30mm, minimum height=9mm, align=center,
        fill=simBgGreen, draw=simGreen] (s) at (5.7,0.7) {$\dot\vs = \vv$};
  \node[stage, minimum width=44mm, minimum height=9mm, align=center,
        fill=simBgGreen, draw=simGreen] (v) at (6.6,-0.8)
    {$\dot\vv = \Mop^{-1}(\vJ - \Cop\vv - \Kop\vs)$};
  \node[font=\scriptsize\itshape, text=simGreen, above=1mm of s]
    {first order in $\vy = (\vs,\vv)$};
  % thread velocity from top into bottom and stiffness from bottom into top
  \draw[refedge] (s.south) -- node[left, font=\scriptsize]{$\vs$} (v.north west);
  \draw[backflow] (v.north east) .. controls +(0,8mm) and +(0,-8mm) ..
    (s.east) node[pos=0.5, right, font=\scriptsize]{$\vv$};
\end{tikzpicture}
\caption[The second-order to first-order recast]{The recast of
\cref{sec:time-firstorder}. The semi-discrete transient system is second order in
$\vs$ (left), so no off-the-shelf integrator---all written for $\dot\vy =
\vF(\vy,t)$---fits it directly. Introducing the velocity $\vv \equiv \dot\vs$ as a
second unknown and stacking $(\vs, \vv)$ into a doubled state turns it into the
first-order pair (right): the top equation is the definition of $\vv$, the bottom
is the original system solved for $\dot\vv$. The doubled state is twice as long;
in exchange every first-order integrator now applies. This recast is the first
thing Palace's \code{TimeDependentFirstOrderOperator} does.}
\label{fig:recast}
\end{figure}

\begin{notationbox}
Throughout this chapter $\vy = (\vs, \vv)$ is the \emph{doubled} first-order state
(position and velocity); $\vF(\vy, t)$ is its rate \cref{eq:time-recast}; $\Delta
t$ is the time step; and $\vy_n \approx \vy(t_n)$ is the numerical state at time
$t_n = t_0 + n\,\Delta t$. We write $\lambda$ for a scalar decay rate when we study
a single mode in isolation. The three operators $\Mop, \Cop, \Kop$ are the
assembled mass, damping, and stiffness of \cref{ch:reductions-pde}, held fixed for
the whole march.
\end{notationbox}

Palace performs exactly this recast: it builds a
\code{TimeDependentFirstOrderOperator} encoding $\vF(\vy, t)$ from the three
assembled matrices and hands it to a library integrator constructed once, outside
the time loop.\footnote{The first-order operator is assembled from $\Mop,\Cop,\Kop$
and the integrator chosen at construction
(\code{palace/models/timeoperator.cpp:311-335}); the loop that advances it lives in
the driver (\code{palace/drivers/transientsolver.cpp:77}). We return to the loop in
\cref{sec:fold-march}.} Everything that follows is about \emph{what one step of
that integrator does}, and how the steps chain together.

\section{Explicit and implicit steps, and why electromagnetics needs implicit}
\label{sec:explicit-implicit}

An integrator advances the state from $\vy_n$ to $\vy_{n+1}$ over one step
$\Delta t$. The simplest possible rule is \emph{forward (explicit) Euler}: read the
rate at the current state and step along it,
\begin{equation}
  \vy_{n+1} = \vy_n + \Delta t\,\vF(\vy_n, t_n).
  \label{eq:fwd-euler}
\end{equation}
Everything on the right is \emph{known}---it is built entirely from $\vy_n$, the
state we already hold---so $\vy_{n+1}$ is computed by a single evaluation of
$\vF$. No system is solved; the step is cheap. This is the defining feature of an
\emph{explicit} method: the new state is an explicit formula in the old one.

The alternative is to evaluate the rate at the \emph{new}, as-yet-unknown state.
\emph{Backward (implicit) Euler} does exactly this:
\begin{equation}
  \vy_{n+1} = \vy_n + \Delta t\,\vF(\vy_{n+1}, t_{n+1}).
  \label{eq:bwd-euler}
\end{equation}
Now $\vy_{n+1}$ appears on \emph{both} sides. The rule no longer hands us the
answer; it defines the answer \emph{implicitly}, as the solution of an equation.
For our linear transient system $\vF$ is affine in $\vy$, so
\cref{eq:bwd-euler} rearranges into a linear system in $\vy_{n+1}$,
\begin{equation}
  \bigl(\mat I - \Delta t\,\mat J\bigr)\,\vy_{n+1} = \vy_n + \Delta t\,\vg_{n+1},
  \label{eq:implicit-solve}
\end{equation}
where $\mat J$ is the (constant) Jacobian of $\vF$ and $\vg_{n+1}$ collects the
forcing. The shape is the headline of the section: \textbf{every implicit step is a
linear solve}. The matrix $\mat I - \Delta t\,\mat J$---built from the same
$\Mop, \Cop, \Kop$ that assembled the operator---must be inverted (in the
matrix-free sense of \cref{ch:krylov}: solved, not literally inverted) once per
step. That solve is a \code{ksp\_solve}, the Krylov solve of
\cref{ch:krylov,ch:cg,ch:gmres}, sitting inside the time loop.

\begin{keyidea}
An \emph{explicit} step evaluates the rate at the state it already knows: it is a
formula, cheap, no solve. An \emph{implicit} step evaluates the rate at the
unknown new state: it is an equation, and solving it for a linear system means a
linear solve every step---a \code{ksp\_solve} (\cref{ch:krylov}) nested inside the
time march. Implicit steps cost more per step; the next subsection shows why
electromagnetics pays the cost anyway.
\end{keyidea}

\subsection{Stiffness: the explicit step's fatal step-size limit}

If implicit steps are so much more expensive, why use them? Because the explicit
step's cheapness is illusory for the systems we integrate. The spatial operator
$\Kop$, assembled from a fine finite-element mesh, has an enormous spread of
eigenvalues: the smallest correspond to the slow physical modes we care about, the
largest to fast modes living on the smallest mesh elements. A system with such a
spread is called \emph{stiff}, and stiffness is poison to explicit methods.

The cleanest way to see why is to watch a single mode. Diagonalize the system and
each mode evolves independently as a scalar decay $\dot y = -\lambda y$, with
$\lambda > 0$ the mode's decay rate; the stiff modes are the ones with very large
$\lambda$. Apply forward Euler \cref{eq:fwd-euler} to this scalar equation:
\begin{equation}
  y_{n+1} = y_n + \Delta t\,(-\lambda y_n) = (1 - \lambda\,\Delta t)\,y_n.
  \label{eq:fe-scalar}
\end{equation}
Each step multiplies $y$ by the \emph{amplification factor} $g = 1 - \lambda\Delta
t$. The true solution $y(t) = y_0 e^{-\lambda t}$ decays, so a faithful numerical
solution must \emph{also} decay: we need $|g| < 1$. That demands
\begin{equation}
  |1 - \lambda\,\Delta t| < 1
  \qquad\Longleftrightarrow\qquad
  \Delta t < \frac{2}{\lambda}.
  \label{eq:fe-stability}
\end{equation}
The step is bounded by $2/\lambda$, and $\lambda$ here is the \emph{largest}
eigenvalue in the system. One fast mode on one tiny element---a mode physically
irrelevant to the answer---pins $\Delta t$ to a minuscule value, and the explicit
march must take an absurd number of tiny steps to cross a physically interesting
interval. This is the stiffness catastrophe.

Backward Euler suffers no such limit. Apply \cref{eq:bwd-euler} to $\dot y =
-\lambda y$:
\begin{equation}
  y_{n+1} = y_n + \Delta t\,(-\lambda y_{n+1})
  \quad\Longrightarrow\quad
  y_{n+1} = \frac{1}{1 + \lambda\,\Delta t}\,y_n,
  \qquad g = \frac{1}{1 + \lambda\,\Delta t}.
  \label{eq:be-scalar}
\end{equation}
For \emph{every} $\lambda > 0$ and \emph{every} $\Delta t > 0$, the denominator
exceeds $1$, so $|g| < 1$: the numerical solution decays no matter how large the
step. The fast modes are damped, the slow modes are tracked, and $\Delta t$ is
chosen for \emph{accuracy on the modes we care about}, not forced small by modes we
do not. This unconditional stability is why electromagnetic transient solvers are
almost always implicit, and it is what \cref{wex:stiff} makes concrete with
numbers.

\begin{keyidea}
The assembled spatial operator is \emph{stiff}: its largest eigenvalue is huge, set
by the smallest mesh element. An explicit step is stable only for $\Delta t <
2/\lambda_{\max}$---a fast, physically irrelevant mode pins the step size to a
crawl. An implicit (backward-Euler-type) step is stable for \emph{any} $\Delta t$,
so the step is chosen for accuracy, not stability. The per-step linear solve buys
freedom from the stiffness limit. That trade---a solve per step in exchange for big
steps---is why electromagnetics integrates implicitly.
\end{keyidea}

\subsection{Stability regions and A-stability}

The scalar test $\dot y = -\lambda y$ generalizes into a standard diagnostic. Let
$\lambda$ range over the \emph{complex} plane (a mode may oscillate as well as
decay, $\lambda = a + ib$), feed $z = \lambda\,\Delta t$ into the method's
amplification factor $g(z)$, and ask: for which $z$ is $|g(z)| \le 1$? That set is
the method's \emph{stability region}---the values of $z$ for which the method does
not amplify. A method is useful for a given step exactly when every mode's $z$
lands inside it.

\Cref{fig:stability} draws the two Euler regions in the $z$-plane. Forward Euler's
region is a single disk of radius $1$ centered at $z = -1$: a bounded set, so
large $\lambda\Delta t$ falls outside it and blows up---the limit
\cref{eq:fe-stability} read geometrically. Backward Euler's region is the
\emph{entire complex plane outside} the disk of radius $1$ centered at $+1$, which
contains the whole left half-plane $\Re(z) < 0$ where the physical decaying modes
live. A method whose stability region contains the entire left half-plane is called
\emph{A-stable}: it is stable for every decaying mode at every step size. Backward
Euler is A-stable; forward Euler is not. A-stability is the precise property the
stiffness argument was groping toward, and it is the property every integrator in
this chapter possesses.

\begin{figure}[t]
\centering
\begin{tikzpicture}[font=\footnotesize]
  % ===== LEFT: forward Euler region (disk |1+z|<=1) =====
  \begin{scope}
    \node[font=\small\bfseries, text=simRust] at (0,2.5) {forward Euler};
    \draw[->, simGray] (-2.6,0) -- (1.4,0) node[right, font=\scriptsize]{$\Re z$};
    \draw[->, simGray] (0,-1.9) -- (0,1.9) node[above, font=\scriptsize]{$\Im z$};
    \fill[simRust!20, draw=simRust, thick] (-1,0) circle (1);
    \node[font=\scriptsize, text=simRust] at (-1,0) {stable};
    \node[font=\scriptsize, text=simGray] at (-1,-1.35) {disk, $|1+z|\le 1$};
    \node[font=\scriptsize, text=simRust!80!black, align=center] at (0.55,1.45)
      {outside:\\blows up};
  \end{scope}
  \draw[simGray!50, thick] (2.0,-2.0) -- (2.0,2.7);
  % ===== RIGHT: backward Euler region (everything outside |1-z|<1) =====
  \begin{scope}[xshift=5.4cm]
    \node[font=\small\bfseries, text=simGreen] at (0,2.5) {backward Euler};
    % shade the whole left half (the important part) lightly green
    \fill[simBgGreen] (-2.6,-1.9) rectangle (0,1.9);
    \draw[->, simGray] (-2.6,0) -- (2.0,0) node[right, font=\scriptsize]{$\Re z$};
    \draw[->, simGray] (0,-1.9) -- (0,1.9) node[above, font=\scriptsize]{$\Im z$};
    % the UNSTABLE disk |1-z|<1 centered at +1
    \fill[white, draw=simGreen, thick, dashed] (1,0) circle (1);
    \node[font=\scriptsize, text=simGray, align=center] at (1.0,0) {un-\\stable};
    \node[font=\scriptsize, text=simGreen, align=center] at (-1.4,1.0)
      {stable\\(all of $\Re z<0$)};
    \node[font=\scriptsize, text=simGreen!60!black] at (-1.3,-1.55) {A-stable};
  \end{scope}
\end{tikzpicture}
\caption[Stability regions of the two Euler methods]{The stability regions in the
$z = \lambda\,\Delta t$ plane (\cref{sec:explicit-implicit}). \emph{Forward Euler}
(left) is stable only on the disk $|1+z|\le 1$---a \emph{bounded} set, so a stiff
mode with large $\lambda\Delta t$ falls outside and the step blows up. \emph{Backward
Euler} (right) is stable \emph{everywhere except} the small disk $|1-z|<1$; its
region contains the entire left half-plane $\Re z<0$ where every physically decaying
mode lives. A method whose region contains all of $\Re z<0$ is \emph{A-stable}.
Backward Euler is A-stable and forward Euler is not---the geometric statement of
why the implicit step has no step-size limit.}
\label{fig:stability}
\end{figure}

\subsection{Accuracy: order, and the truncation error}

Stability says the march does not blow up; \emph{accuracy} says it lands near the
true trajectory. The standard measure is the \emph{order} $p$ of a method: the
local truncation error---the error committed in a single step from exact data---is
$O(\Delta t^{p+1})$, and the global error accumulated over a fixed interval is
$O(\Delta t^{p})$. A first-order method ($p=1$, like both Eulers) halves its global
error when you halve the step; a second-order method ($p=2$) quarters it. Higher
order buys accuracy faster, at the cost of more work or more storage per step.

Order and stability are \emph{independent} axes, and a good integrator wants both:
high order for accuracy and A-stability for the stiff transient system. The plain
Eulers give one or the other---backward Euler is A-stable but only first order;
nothing here is high order. The integrators Palace actually uses, next, are the
ones engineered to be \emph{both} A-stable \emph{and} second order, with one extra
feature the plain methods lack.

\section{The integrators Palace uses}
\label{sec:integrators}

Palace offers a small menu of implicit integrators, selected at construction. We
present the \emph{idea} of each---what property it is engineered for---without
reproducing the Butcher tableaux that specify their coefficients; the tableaux are
in any numerical-methods reference, and the ideas are what matter for reading the
spec.

\subsection{Generalized-$\alpha$: tunable numerical damping}

The workhorse is the \emph{generalized-$\alpha$} method. It is second-order
accurate and A-stable---the two properties the plain Eulers could not combine---and
it adds a third, decisive feature: \emph{tunable, frequency-selective numerical
damping}. A pure second-order A-stable method (the trapezoidal rule) tracks every
mode without damping, which sounds ideal but is not: the highest-frequency modes a
mesh can represent are numerical junk---unresolved, aliased, physically
meaningless---and an undamped method lets them oscillate forever, polluting the
solution. Generalized-$\alpha$ damps them on purpose.

The point is best seen in its \emph{spectral-radius} curve, \cref{fig:genalpha}.
The method exposes a single dial, a high-frequency spectral radius $\rho_\infty \in
[0,1]$, that sets how aggressively the infinitely-fast modes are damped while
leaving the low, resolved modes essentially untouched. Set $\rho_\infty = 1$ and no
mode is damped (the trapezoidal limit); set $\rho_\infty = 0$ and the
fastest modes are annihilated in one step (the maximally dissipative limit); in
between, the low modes pass with amplification near $1$ and the high modes are
selectively bled away. The dial lets the analyst keep the physics and discard the
numerical noise, with one number, without touching accuracy on the resolved
band.\footnote{Palace constructs a \code{GeneralizedAlphaSolver} (alongside
\code{SDIRK23Solver} and the SUNDIALS \code{ARKStepSolver}) at
\code{palace/models/timeoperator.cpp:311-335}; the choice is a configuration token,
the integrator library-owned.}

\begin{figure}[t]
\centering
\begin{tikzpicture}
\begin{axis}[
    width=0.74\linewidth, height=52mm,
    xlabel={mode frequency $\;\omega\,\Delta t$ (low $\to$ high)},
    ylabel={amplification $|g|$},
    xmin=0, xmax=40, ymin=0, ymax=1.10,
    axis lines=left, ytick={0,0.5,1},
    xtick={0,10,20,30,40},
    yticklabel style={font=\scriptsize}, xticklabel style={font=\scriptsize},
    xlabel style={font=\scriptsize}, ylabel style={font=\scriptsize},
    legend style={font=\scriptsize, draw=simGray!50, at={(0.98,0.97)},
      anchor=north east, legend cell align=left},
    every axis plot/.append style={very thick, samples=120, domain=0:40},
  ]
  % rho_inf = 1: trapezoidal, no damping -> |g| stays 1
  \addplot[simTeal] {1};
  \addlegendentry{$\rho_\infty=1$: no damping (trapezoidal)}
  % rho_inf = 0.5: high modes damped toward 0.5
  \addplot[simBlue] {0.5 + 0.5/(1 + (x/5)^2)};
  \addlegendentry{$\rho_\infty=0.5$: moderate damping}
  % rho_inf = 0: high modes annihilated
  \addplot[simRust] {1/(1 + (x/5)^2)};
  \addlegendentry{$\rho_\infty=0$: maximal damping}
  % mark the resolved band
  \draw[simGray, dashed] (axis cs:7,0) -- (axis cs:7,1.10);
  \node[font=\scriptsize, text=simGray, anchor=south west]
    at (axis cs:7.5,0.05) {resolved $\to$ unresolved};
\end{axis}
\end{tikzpicture}
\caption[Numerical damping of generalized-$\alpha$]{The frequency-selective damping
of generalized-$\alpha$ (\cref{sec:integrators}), drawn as the per-step
amplification $|g|$ against mode frequency. The single dial $\rho_\infty$ sets the
high-frequency behavior: at $\rho_\infty = 1$ (teal) no mode is damped, the
undamped trapezoidal limit; at $\rho_\infty = 0$ (rust) the fastest modes are
annihilated in one step; intermediate values (blue) bleed the high modes away while
leaving the low, \emph{resolved} modes (left of the dashed line) essentially
untouched at $|g|\approx 1$. This is the feature the plain Eulers lack: the analyst
keeps the resolved physics and discards the unresolved numerical noise with one
number.}
\label{fig:genalpha}
\end{figure}

\subsection{SDIRK and ARKStep}

The other two menu items are Runge--Kutta-family methods. \emph{SDIRK}---singly
diagonally implicit Runge--Kutta---takes a small number of internal stages per
step, each stage an implicit solve, but arranges the stages so they share a
\emph{single} matrix factorization (the ``singly diagonal'' structure means every
stage solves a system with the same $\mat I - \gamma\Delta t\,\mat J$, so the
preconditioner is built once and reused across stages). SDIRK methods reach second
order and beyond while staying A-stable, at the cost of several solves per step
rather than one---the price of higher accuracy bought stage by stage.
\emph{ARKStep} is the additive Runge--Kutta integrator from the SUNDIALS library,
designed to split a system into a stiff part it treats implicitly and a non-stiff
part it treats explicitly---an implicit-explicit (IMEX) method---and to adapt its
step size automatically to a requested error tolerance.

For our purposes the three integrators share the structure that matters and differ
only in details we do not need:

\begin{description}[style=nextline, leftmargin=0pt, font=\normalfont\itshape]
  \item[All are implicit.] Every one solves at least one linear system per step (or
  per stage)---a \code{ksp\_solve} nested in the march. None escapes the per-step
  solve, because none escapes A-stability, and A-stability for these methods
  requires the implicit evaluation.
  \item[All are A-stable.] Each is stable on the whole stiff transient system at a
  step chosen for accuracy, not pinned by the largest eigenvalue. This is the
  non-negotiable property; it is why the menu has no explicit entry.
  \item[They differ in order, stage count, and damping control.]
  Generalized-$\alpha$ offers the damping dial; SDIRK trades stages for order;
  ARKStep adapts the step and splits stiff from non-stiff. These are knobs on a
  shared frame, not different algorithms.
\end{description}

\begin{notationbox}
The key fact to carry forward is structural, not numerical: \emph{whichever}
integrator is chosen, one step is an opaque routine that consumes the current
state and the step size and returns the next state, performing one or more
\code{ksp\_solve}s internally. The framework does not look inside that routine---it
quantifies over it. That opacity is exactly what lets the next section treat
``one step'' as a single black box and focus on how the steps \emph{chain}.
\end{notationbox}

\section{The march is a state-threaded fold}
\label{sec:fold-march}

We now reach the chapter's structural payoff, the reason the transient analysis was
promised as the canonical example back in \cref{ch:fold} and given a combinator
name in \cref{ch:combinators}. Strip away the integrator's internals---treat one
step as an opaque operator \code{time\_step\_op} that advances the state---and look
at how the march is assembled from steps. It is a \emph{fold}.

\subsection{Why a fold and not a map}

Recall the fundamental split of \cref{ch:combinators}. A \code{map} runs a family
of \emph{independent} computations: no carry between elements, so they reorder and
parallelize. A \code{fold} threads a \emph{carry} through a chain: each step
consumes what the last produced, so the steps are strictly ordered. The transient
march is unambiguously the second. The field-state at $t_{n+1}$ is computed
\emph{from} the field-state at $t_n$ by one time step; step $n+1$ cannot begin
until step $n$ has finished, because step $n$'s output is step $n+1$'s input. The
state is threaded forward along a single timeline. This is the
\emph{state-threaded fold}, the combinator \code{fold\_solve}.

Contrast it, as \cref{ch:combinators} did, with the \code{solve\_family} of
electrostatics. There the several solves shared one operator and depended on
nothing but their own right-hand side---independent, a map, embarrassingly
parallel. Here the steps share one operator \emph{and} a threaded state, and the
state makes them sequential. The presence or absence of the carry is the entire
difference, and it is visible in the type. \Cref{fig:fold-vs-map} sets the two side
by side on the same timeline: the map fans out, the fold lines up.

\begin{figure}[t]
\centering
\begin{tikzpicture}[font=\footnotesize]
  % ===== LEFT: the transient FOLD (carry threaded) =====
  \begin{scope}
    \node[font=\small\bfseries, text=simGreen] at (2.4,2.5) {transient: a \code{fold} (carry threaded)};
    \node[data, minimum width=9mm] (s0) at (-0.3,1.1) {$\vy_0$};
    \foreach \i/\x in {1/1.1,2/2.6,3/4.1}{
      \node[opbox, minimum width=11mm] (g\i) at (\x,1.1) {step};
    }
    \node[data, minimum width=9mm] (s3) at (5.4,1.1) {$\vy_3$};
    \draw[flow] (s0.east) -- (g1.west);
    \draw[flow] (g1.east) -- node[above,font=\scriptsize]{$\vy_1$} (g2.west);
    \draw[flow] (g2.east) -- node[above,font=\scriptsize]{$\vy_2$} (g3.west);
    \draw[flow] (g3.east) -- (s3.west);
    % each step hides a ksp_solve
    \foreach \i/\x in {1/1.1,2/2.6,3/4.1}{
      \node[extern, minimum width=11mm, minimum height=5mm, font=\scriptsize]
        (k\i) at (\x,-0.1) {\code{ksp}};
      \draw[refedge] (g\i.south) -- (k\i.north);
    }
    \node[font=\scriptsize, text=simRust, align=center] at (2.4,-0.95)
      {each step reads the prior $\vy$ \emph{and} hides an implicit solve:\\
       strictly ordered---cannot overlap timesteps};
  \end{scope}
  \draw[simGray!50, thick] (6.2,-1.3) -- (6.2,2.7);
  % ===== RIGHT: the electrostatic MAP (no carry) =====
  \begin{scope}[xshift=7.0cm]
    \node[font=\small\bfseries, text=simBlue] at (2.0,2.5) {electrostatic: a \code{map} (no carry)};
    \node[data, minimum width=11mm] (fam) at (2.0,1.6) {one $\Kop$};
    \foreach \i/\x in {1/0.4,2/2.0,3/3.6}{
      \node[opbox, minimum width=10mm] (m\i) at (\x,0.6) {solve};
      \node[product, minimum width=10mm, minimum height=5mm] (r\i) at (\x,-0.5) {$\vx_{\i}$};
      \draw[flow] (fam.south) -- (m\i.north);
      \draw[flow] (m\i.south) -- (r\i.north);
    }
    \node[font=\scriptsize, text=simGreen, align=center] at (2.0,-1.15)
      {no arrows \emph{between} solves:\\independent, reorderable, parallel};
  \end{scope}
\end{tikzpicture}
\caption[The transient fold versus the electrostatic map]{The transient march
(left) as a state-threaded \code{fold}, beside the electrostatic family (right) as
an independent \code{map}. In the fold the doubled state $\vy_0 \to \vy_1 \to \vy_2
\to \vy_3$ is threaded forward---each step consumes the prior state, so the steps
are strictly ordered and \emph{cannot} overlap---and each step additionally hides
an implicit \code{ksp\_solve} (\cref{sec:explicit-implicit}). In the map the
several solves share one operator $\Kop$ but carry nothing between them, so they
fan out and parallelize. The horizontal carry arrows, present on the left and
absent on the right, are the whole difference: \code{fold\_solve} versus
\code{solve\_family}.}
\label{fig:fold-vs-map}
\end{figure}

\subsection{\texttt{fold\_solve}, in the calculus}

In the vocabulary of \cref{ch:combinators}, the march is one expression. Capture
the operator once, seed the state once, and \code{foldl} the opaque per-step
operator over the schedule of times:
\begin{lstlisting}[style=l4style]
fold_solve :: OpParams -> TimeState -> [Time] -> TimeState
fold_solve op s0 schedule = foldl (\s t -> time_step_op op s t) s0 schedule

-- the opaque per-step operator the fold quantifies over:
-- advances the field-state one step; bottoms out in the library integrator
time_step_op :: OpParams -> TimeState -> Time -> TimeState
\end{lstlisting}
Read the signature and the whole structure of the march is there. The seed
\code{s0} and the result are both a \code{TimeState}---the field-state carried
forward. The step \lstinline[style=l4style]|time_step_op op s t| reads \code{s},
the prior carry; \emph{that read} is what makes this a fold and forbids reordering
at the type level. The operator \code{op} (the captured $\Mop,\Cop,\Kop$ and the
chosen integrator) is bound once, before the fold, and threaded unchanged into
every step---the construction hoist of \cref{sec:integrators} made structural.
And \code{time\_step\_op} is \emph{opaque}: it bottoms out in the library
integrator step of \cref{sec:integrators}, which the combinator only quantifies
over. \code{fold\_solve} owns the loop; the library owns one step.

This is one of the four solve corners of \cref{ch:situating,ch:combinators}---the
\emph{state-threaded fold} corner---and the same combinator drives the
adaptive-mesh-refinement loop of \cref{ch:amr}, where each refinement pass begins
from the mesh the last pass produced. The transient march is simply its purest
instance.

\begin{keyidea}
The time-march is a \emph{sequential fold}: each step begins where the last ended.
It is the \code{fold\_solve} combinator---\lstinline[style=l4style]|foldl time_step_op s0 schedule|---threading
a single field-state carry forward along one timeline, with the operator captured
once outside the loop and each opaque step hiding an implicit \code{ksp\_solve}.
This is the genuinely sequential corner of the simulator: unlike the electrostatic
map, the transient fold cannot overlap its timesteps, because step $n+1$'s input is
step $n$'s output. \emph{Implicit steps cost a linear solve; the carry makes the
march serial.}
\end{keyidea}

Two laws follow immediately from \code{fold\_solve} being a left fold, and they are
the opposite of the map's laws. The map of electrostatics \emph{distributes} over
concatenation, $f(a \mathbin{+\!\!+} b) = f\,a \mathbin{+\!\!+} f\,b$---the algebraic
license for splitting the family across machines. The fold instead
\emph{threads} through concatenation, which is the subject of the next section and
the cleanest gift the fold structure gives us. First, though, note the
\emph{non}-laws: the schedule does not commute (permuting the timesteps changes the
trajectory) and the steps do not fuse (each is an opaque integrator step with its
own internal state advance). A fold is not a map wearing a disguise; the carry is
real, and the type records it.

\section{Checkpoint and resume: the schedule-split law}
\label{sec:checkpoint}

The single most useful algebraic fact about \code{fold\_solve} is the law that a
left fold satisfies over a concatenated list. Split the schedule into a prefix
$a$ and a suffix $b$; then folding the whole is the same as folding the prefix and
folding the suffix \emph{from the state the prefix produced}:
\begin{equation}
  \boxed{\;\code{fold\_solve}\ \code{op}\ \vy_0\ (a \mathbin{+\!\!+} b)
  \;=\;
  \code{fold\_solve}\ \code{op}\ \bigl(\code{fold\_solve}\ \code{op}\ \vy_0\ a\bigr)\ b\;}
  \label{eq:schedule-split}
\end{equation}
This is the standard identity $\code{foldl}\,(a \mathbin{+\!\!+} b) = \code{foldl}\,b
\circ \code{foldl}\,a$, specialized to the captured-operator march. Read it
operationally and it says something an engineer cares about a great deal:
\textbf{a march can be checkpointed and resumed}. Run the prefix $a$, save the
final state $\vy_{|a|} = \code{fold\_solve}\ \code{op}\ \vy_0\ a$ to disk, and stop.
Later, load $\vy_{|a|}$ and run the suffix $b$ seeded from it. The law guarantees
the result is \emph{bit-for-bit} the same trajectory as the uninterrupted run---the
split changed nothing, because the only thing that crosses the split point is the
carry, and the carry is exactly what we saved.

\Cref{fig:checkpoint} draws the equality. The top row is the uninterrupted march;
the bottom row splits it, saving the carry at the cut and resuming from it. The two
produce the same final state because the carry threaded across the cut is the same
value either way.

\begin{figure}[t]
\centering
\begin{tikzpicture}[font=\footnotesize]
  % ===== TOP: uninterrupted march =====
  \node[font=\small\bfseries, text=simInk, anchor=west] at (-0.3,2.5)
    {one uninterrupted march};
  \node[data, minimum width=8mm] (u0) at (0,1.7) {$\vy_0$};
  \foreach \i/\x in {1/1.0,2/2.2,3/3.4,4/4.6}{
    \node[opbox, minimum width=8mm, minimum height=5mm] (us\i) at (\x,1.7) {};
  }
  \node[data, minimum width=8mm] (u4) at (5.6,1.7) {$\vy_4$};
  \draw[flow] (u0)--(us1); \draw[flow] (us1)--(us2); \draw[flow] (us2)--(us3);
  \draw[flow] (us3)--(us4); \draw[flow] (us4)--(u4);
  \node[font=\scriptsize, text=simGray] at (2.8,1.15) {schedule $a \mathbin{+\!\!+} b$};
  % ===== BOTTOM: split + checkpoint =====
  \node[font=\small\bfseries, text=simGreen, anchor=west] at (-0.3,0.3)
    {split: run $a$, save carry, resume $b$};
  \node[data, minimum width=8mm] (b0) at (0,-0.5) {$\vy_0$};
  \foreach \i/\x in {1/1.0,2/2.2}{
    \node[opbox, minimum width=8mm, minimum height=5mm] (bs\i) at (\x,-0.5) {};
  }
  % the checkpoint
  \node[product, minimum width=10mm, fill=simBgGold, draw=simGold] (ckpt) at (3.4,-0.5) {$\vy_2$};
  \foreach \i/\x in {3/4.6,4/5.6}{
    \node[opbox, minimum width=8mm, minimum height=5mm] (bs\i) at (\x,-0.5) {};
  }
  \node[data, minimum width=8mm] (b4) at (6.6,-0.5) {$\vy_4$};
  \draw[flow] (b0)--(bs1); \draw[flow] (bs1)--(bs2); \draw[flow] (bs2)--(ckpt);
  \draw[flow] (ckpt)--(bs3); \draw[flow] (bs3)--(bs4); \draw[flow] (bs4)--(b4);
  \node[font=\scriptsize, text=simGreen] at (1.6,-1.05) {fold $a$};
  \node[font=\scriptsize, text=simRust] at (5.1,-1.05) {fold $b$ from $\vy_2$};
  % save-to-disk callout
  \node[font=\scriptsize, text=simGold!75!black, align=center] at (3.4,-1.5)
    {save / load\\(checkpoint)};
  % equality brace
  \node[font=\Large, text=simGray] at (6.3,0.6) {$=$};
\end{tikzpicture}
\caption[Checkpoint and resume as a schedule split]{The schedule-split law
\cref{eq:schedule-split} drawn as checkpoint-and-resume. \emph{Top}: an
uninterrupted march folds the whole schedule $a \mathbin{+\!\!+} b$ from $\vy_0$ to
$\vy_4$. \emph{Bottom}: the same march split at the midpoint---fold the prefix $a$
to the carry $\vy_2$, \emph{save it} (the gold checkpoint), stop, then later load
$\vy_2$ and fold the suffix $b$ from it. The two produce the identical $\vy_4$,
because the only thing crossing the cut is the carry, and the carry is exactly what
is saved. This is why long runs can be checkpointed against faults and resumed
without changing the answer: it is a clean algebraic property of the fold, not an
engineering approximation.}
\label{fig:checkpoint}
\end{figure}

\begin{keyidea}
Splitting the time schedule and resuming from the saved state yields \emph{exactly}
the same trajectory: $\code{fold\_solve}\ \code{op}\ \vy_0\ (a\mathbin{+\!\!+}b) =
\code{fold\_solve}\ \code{op}\ (\code{fold\_solve}\ \code{op}\ \vy_0\ a)\ b$. This
is the fold's thread-through law---the carry is the only thing that crosses a split
point, so saving the carry \emph{is} a checkpoint. Long runs survive faults
(checkpoint, restart from disk), expensive runs split across job-scheduler windows,
and none of it perturbs the result. Restartability is not bolted on; it is what the
fold structure means.
\end{keyidea}

This is the practical reason the fold's sequentiality, which sounded like pure cost
in \cref{sec:fold-march}, comes with a genuine compensation. A map parallelizes but
has no notion of ``resume''---there is no carry to save. A fold serializes but is
\emph{restartable}, because the carry it threads is precisely the minimal state
needed to continue. The same property---a single carry threaded forward---that
forbids parallelism across timesteps is what makes a checkpoint a single saved
value rather than a replay of the whole history.

\section{Reading the results}
\label{sec:time-readout}

The fold produces a \emph{trajectory}: the field-state at each step. From it the
transient driver reads the physical products, and here the demand-driven pruning of
\cref{ch:fold} earns its keep. If a consumer reads only the final state, the
intermediate states are pruned and the fold degenerates to a bare carry-threading
loop; if a consumer reads the per-step outputs---and the transient driver does, at
every step---the trajectory materializes. What the driver reads at each step is:

\begin{description}[style=nextline, leftmargin=0pt, font=\normalfont\itshape]
  \item[The field trajectory.] From each step's state the electric and magnetic
  fields $\vE, \vB$ are recovered (the \code{GetE}/\code{GetB} of the time
  operator) and written out---the ``movie'' the transient analysis exists to
  produce, a frame per step.
  \item[Port voltages and currents.] Integrating the fields against the port modes
  gives a time-domain voltage and current at each terminal: the waveforms an
  engineer reads off an oscilloscope, computed step by step.
  \item[Energy versus time.] The field energy
  $\tfrac12(\vs^{\!\top}\Mop\,\vs + \dots)$, evaluated at each step, traces how
  energy enters with the pulse, sloshes between electric and magnetic forms, and
  bleeds away through loss and radiation. \Cref{fig:energy} draws a representative
  trajectory: an injected pulse rings up the stored energy, which then decays as
  the damping $\Cop$ and the ports carry it off.
\end{description}

\begin{figure}[t]
\centering
\begin{tikzpicture}
\begin{axis}[
    width=0.80\linewidth, height=50mm,
    xlabel={time $t$}, ylabel={},
    xmin=0, xmax=10, ymin=-1.15, ymax=1.25,
    axis lines=left, xtick={0,2,4,6,8,10}, ytick={-1,0,1},
    yticklabel style={font=\scriptsize}, xticklabel style={font=\scriptsize},
    xlabel style={font=\scriptsize},
    legend style={font=\scriptsize, draw=simGray!50, at={(0.98,0.97)},
      anchor=north east, legend cell align=left},
    every axis plot/.append style={very thick, samples=120, domain=0:10},
  ]
  % injected pulse: a modulated gaussian
  \addplot[simRust] {0.9*exp(-2*(x-1.5)^2)*cos(deg(6*(x-1.5)))};
  \addlegendentry{drive $\vJ(t)$ (a pulse)}
  % a field component: rings up then decays
  \addplot[simBlue] {0.8*(1-exp(-3*x))*exp(-0.35*x)*cos(deg(4*x-1))};
  \addlegendentry{field $E(t)$ (rings, decays)}
  % energy: rises then decays monotonically
  \addplot[simGreen, densely dashed] {0.95*(1-exp(-2.2*x))*exp(-0.5*x)};
  \addlegendentry{stored energy (rise then decay)}
\end{axis}
\end{tikzpicture}
\caption[A transient field and energy trajectory]{A representative readout of the
transient march (\cref{sec:time-readout}). The pulse drive $\vJ(t)$ (rust) injects
a short burst; a field component $E(t)$ (blue) rings up as the pulse arrives and
then decays as energy leaves the system; the stored energy (green, dashed) rises
with the injection and decays monotonically thereafter, bled away by the damping
$\Cop$ and carried off through the ports. Each curve is sampled once per fold step,
so the resolution of the readout is exactly the resolution of the march---one frame
per timestep.}
\label{fig:energy}
\end{figure}

\subsection{A bridge to the frequency domain}

One readout deserves a forward pointer. A transient run injects a \emph{broadband}
pulse---a superposition of many frequencies at once (\cref{sec:transient})---and
records the time-domain response. Taking the Fourier transform of an input
waveform and the corresponding output waveform, and dividing, recovers the
frequency response across the whole band of the pulse \emph{in a single run}. This
is the time-domain route to the same S-parameters the driven analysis
(\cref{ch:reductions}) computes frequency by frequency: where the driven analysis
sweeps a \code{frequency\_sweep} map, one frequency per solve, the transient
analysis folds one march and Fourier-transforms the result. The two analyses
compute the same physics by the two routes of the harmonic rule
(\cref{ch:reductions-pde})---a single frequency at a time, or all frequencies at
once through a pulse. The transient driver and its products are studied in full in
\cref{ch:transient}.

\section{Worked examples}

\begin{workedexample}{Stiffness made concrete: explicit blows up, implicit
survives}{stiff}
We make the stiffness argument of \cref{sec:explicit-implicit} numerical on the
scalar test equation $\dot y = -\lambda y$ with a \emph{stiff} rate $\lambda = 100$
and a modest step $\Delta t = 0.05$. The exact solution from $y_0 = 1$ is
$y(t) = e^{-100 t}$, which decays to nothing almost instantly.

\medskip
\emph{Explicit Euler.} The amplification factor \cref{eq:fe-scalar} is
\[
  g_{\text{exp}} = 1 - \lambda\,\Delta t = 1 - 100\cdot 0.05 = 1 - 5 = -4.
\]
Each step multiplies $y$ by $-4$. The stability bound \cref{eq:fe-stability} demands
$\Delta t < 2/\lambda = 0.02$, and our step $0.05$ violates it, so $|g| = 4 > 1$ and
the numerical solution \emph{grows} while the true one decays:
\[
  y_0 = 1,\quad y_1 = -4,\quad y_2 = 16,\quad y_3 = -64,\quad y_4 = 256,\ \dots
\]
an explosive oscillation, the exact opposite of $e^{-100t}\to 0$. To make explicit
Euler stable we would have to shrink $\Delta t$ below $0.02$---and on a real mesh
$\lambda_{\max}$ is far larger than $100$, forcing $\Delta t$ smaller still.

\medskip
\emph{Implicit (backward) Euler.} The amplification factor \cref{eq:be-scalar} is
\[
  g_{\text{imp}} = \frac{1}{1 + \lambda\,\Delta t} = \frac{1}{1 + 5} = \frac{1}{6}
  \approx 0.167.
\]
Now $|g| < 1$ for our step---indeed for \emph{any} step---so the solution decays as
it should:
\[
  y_0 = 1,\quad y_1 = \tfrac16 \approx 0.167,\quad y_2 = \tfrac1{36} \approx 0.028,
  \quad y_3 \approx 0.0046,\ \dots
\]
monotone toward zero, tracking the true decay. \Cref{fig:stiff} plots the two
against the exact solution: explicit Euler oscillates out of the frame, implicit
Euler decays cleanly. The implicit step paid for itself with one solve (here a
trivial scalar division; on the full system, a \code{ksp\_solve}) and in exchange
removed the step-size limit entirely. This single picture is the whole reason
electromagnetic transient solvers are implicit.
\end{workedexample}

\begin{figure}[t]
\centering
\begin{tikzpicture}
\begin{axis}[
    width=0.78\linewidth, height=52mm,
    xlabel={step $n$ \;($t = 0.05\,n$)}, ylabel={$y_n$},
    xmin=0, xmax=6, ymin=-9, ymax=9,
    axis lines=middle, xtick={1,2,3,4,5}, ytick={-8,-4,0,4,8},
    yticklabel style={font=\scriptsize}, xticklabel style={font=\scriptsize},
    xlabel style={font=\scriptsize, at={(axis description cs:0.5,-0.06)}},
    ylabel style={font=\scriptsize},
    legend style={font=\scriptsize, draw=simGray!50, at={(0.5,-0.22)},
      anchor=north, legend columns=3, column sep=5pt},
    every axis plot/.append style={very thick, mark=*, mark size=1.5pt},
  ]
  % explicit Euler: 1, -4, 16, -64(clip), ... oscillating blowup
  \addplot[simRust] coordinates {(0,1) (1,-4) (2,8.8) (3,-8.8)};
  \addlegendentry{explicit Euler ($g=-4$)}
  % implicit Euler: 1, 1/6, 1/36, ... decay
  \addplot[simBlue] coordinates {(0,1) (1,0.167) (2,0.028) (3,0.0046) (4,0.0008) (5,0.0001)};
  \addlegendentry{implicit Euler ($g=\tfrac16$)}
  % exact: e^{-100 t} = e^{-5 n}
  \addplot[simGreen, dashed, mark=none, samples=60, domain=0:6] {exp(-5*x)};
  \addlegendentry{exact $e^{-100t}$}
\end{axis}
\end{tikzpicture}
\caption[Explicit blow-up versus implicit stability]{The stiff test of
\cref{wex:stiff}: $\dot y = -\lambda y$, $\lambda = 100$, $\Delta t = 0.05$
(violating the explicit bound $\Delta t < 0.02$). \emph{Explicit Euler} (rust)
amplifies by $g = -4$ each step and oscillates explosively out of the frame---the
later points $\pm 64, \pm 256$ are clipped. \emph{Implicit Euler} (blue) contracts
by $g = \tfrac16$ each step and decays cleanly onto the exact solution (green,
dashed). The implicit step is unconditionally stable; the explicit step is useless
here without a step ten times smaller. This is the stiffness catastrophe, and its
cure, in one plot.}
\label{fig:stiff}
\end{figure}

\begin{workedexample}{Marching a two-degree-of-freedom oscillator}{twodof}
We now run the real machinery---the recast, the per-step implicit solve, and the
state threading---on the smallest non-trivial system: a two-degree-of-freedom
mass--spring--damper, $\Mop\ddot\vs + \Cop\dot\vs + \Kop\vs = \vJ(t)$ with
\[
  \Mop = \begin{psmallmatrix}1 & 0\\ 0 & 1\end{psmallmatrix},\quad
  \Cop = \begin{psmallmatrix}0.2 & 0\\ 0 & 0.2\end{psmallmatrix},\quad
  \Kop = \begin{psmallmatrix}2 & -1\\ -1 & 2\end{psmallmatrix},\quad
  \vJ(t) = \begin{psmallmatrix}1\\ 0\end{psmallmatrix}\ (\text{constant}),
\]
starting from rest, $\vs_0 = \vv_0 = (0,0)$. This is a hand-sized stand-in for the
semi-discrete transient system; $\Kop$ is the stiffness coupling the two degrees of
freedom.

\medskip
\emph{Recast.} The doubled state is $\vy = (\vs, \vv) \in \mathbb R^4$, with rate
$\vF(\vy) = (\vv,\ \Mop^{-1}(\vJ - \Cop\vv - \Kop\vs))$ as in
\cref{eq:time-recast}. Since $\Mop = \mat I$ here, $\dot\vv = \vJ - \Cop\vv -
\Kop\vs$.

\medskip
\emph{One implicit step.} Take backward Euler with $\Delta t = 0.1$. The step
\cref{eq:bwd-euler} couples the two halves: $\vs_{n+1} = \vs_n + \Delta t\,\vv_{n+1}$
and $\vv_{n+1} = \vv_n + \Delta t\,(\vJ - \Cop\vv_{n+1} - \Kop\vs_{n+1})$.
Substituting the first into the second eliminates $\vs_{n+1}$ and leaves a
\emph{linear system} for $\vv_{n+1}$ alone:
\[
  \underbrace{\bigl(\mat I + \Delta t\,\Cop + \Delta t^2\,\Kop\bigr)}_{\displaystyle \mat A}
  \,\vv_{n+1}
  = \vv_n + \Delta t\,\bigl(\vJ - \Kop\,\vs_n\bigr).
\]
This is the implicit step's linear solve, the $2\times 2$ stand-in for the
\code{ksp\_solve} (\cref{sec:explicit-implicit}). With our numbers,
\[
  \mat A = \mat I + 0.1\,\Cop + 0.01\,\Kop
  = \begin{psmallmatrix}1.04 & -0.01\\ -0.01 & 1.04\end{psmallmatrix}.
\]
The matrix $\mat A$ is \emph{fixed for the whole march}---it is built from the
captured $\Mop,\Cop,\Kop$ and $\Delta t$, none of which change---so its
factorization is computed once and reused every step. This is the
operator-capture-once hoist of \code{fold\_solve} (\cref{sec:fold-march}) in
miniature: \code{op} is bound before the fold; every step reads the same $\mat A$.

\medskip
\emph{Thread the state forward.} Starting from $\vy_0 = \mathbf 0$, solve for
$\vv_1$ then set $\vs_1 = \Delta t\,\vv_1$, feed $\vy_1 = (\vs_1, \vv_1)$ into the
next step, and repeat. Three steps:
\begin{center}
\setlength{\tabcolsep}{7pt}
\renewcommand{\arraystretch}{1.2}
\begin{tabular}{@{}clll@{}}
\toprule
step $n$ & RHS $\vv_n + \Delta t(\vJ - \Kop\vs_n)$ & solve $\vv_{n+1} = \mat A^{-1}(\cdot)$ & $\vs_{n+1} = \vs_n + \Delta t\,\vv_{n+1}$ \\
\midrule
0 & $(0.100,\ 0.000)$ & $(0.0961,\ 0.0009)$ & $(0.00961,\ 0.00009)$ \\
1 & $(0.1854,\ 0.0105)$ & $(0.1784,\ 0.0118)$ & $(0.02745,\ 0.00127)$ \\
2 & $(0.2589,\ 0.0264)$ & $(0.2492,\ 0.0278)$ & $(0.05237,\ 0.00405)$ \\
\bottomrule
\end{tabular}
\end{center}
Each row's input is the previous row's output---the carry $\vy_n \to \vy_{n+1}$
threaded forward, exactly the fold of \cref{sec:fold-march}. The driven coordinate
$s^{(1)}$ climbs steadily toward its static deflection while the coupled coordinate
$s^{(2)}$ follows through the off-diagonal stiffness; with the damping $\Cop$
present, the march settles rather than oscillating forever. Three points are the
whole chapter in miniature: \emph{(i)} the recast doubled $\mathbb R^2$ into $\mathbb
R^4$; \emph{(ii)} each step \emph{solved a linear system} $\mat A\vv_{n+1} =
\text{RHS}$, with $\mat A$ fixed and factored once; \emph{(iii)} the state was
\emph{threaded} from each step into the next---a \code{fold\_solve}, not a map. Were
we to checkpoint after step~1, saving $\vy_1$, and resume, the schedule-split law
\cref{eq:schedule-split} guarantees we would land on the identical $\vy_3$.
\end{workedexample}

\summarysection
The transient analysis marches the semi-discrete second-order system $\Mop\ddot\vs
+ \Cop\dot\vs + \Kop\vs = \vJ(t)$ forward in time. First it is recast as a
first-order system $\dot\vy = \vF(\vy, t)$ in the doubled state $\vy = (\vs,
\dot\vs)$, so standard integrators apply. Because the assembled spatial operator is
\emph{stiff}, explicit steps are limited to a tiny $\Delta t < 2/\lambda_{\max}$,
so electromagnetics integrates \emph{implicitly}: each implicit step is a linear
solve---a \code{ksp\_solve} (\cref{ch:krylov})---and the methods used
(generalized-$\alpha$ with its tunable numerical damping, SDIRK, ARKStep) are all
A-stable, trading a per-step solve for freedom from the stiffness limit.
Structurally, the whole march is the \code{fold\_solve} combinator of
\cref{ch:combinators}: a state-threaded \code{foldl} of an opaque per-step operator,
each step beginning where the last ended, the operator captured once outside the
loop. The fold's thread-through law,
$\code{fold\_solve}\,\code{op}\,\vy_0\,(a\mathbin{+\!\!+}b) =
\code{fold\_solve}\,\code{op}\,(\code{fold\_solve}\,\code{op}\,\vy_0\,a)\,b$, makes
checkpoint-and-restart exact: the carry is the only thing crossing a split point.
The march is genuinely sequential---it cannot overlap its timesteps, the price of
the carry the electrostatic map does not pay.

\furtherreading
The recast, stability regions, A-stability, order, and the stiffness phenomenon are
standard material in any numerical-ODE or matrix-computations
reference; Golub and Van Loan~\citep{golubvanloan2013} cover the linear-algebra
machinery of the per-step solve and the eigenvalue spread that makes a system
stiff. Jin's finite-element electromagnetics text~\citep{jin2014} treats the
time-domain finite-element formulation and the specific integrators used for
Maxwell's equations, including the generalized-$\alpha$ family and its numerical
damping. The per-step linear solve is a Krylov solve in the sense of
\cref{ch:krylov}, for which Saad~\citep{saad2003} is the standard reference; the
fold structure and its laws are developed in \cref{ch:fold,ch:combinators}, and the
transient driver that wraps the march is \cref{ch:transient}.

\exercisessection
\begin{enumerate}[nosep]
  \item \emph{The recast preserves the system.} Starting from the doubled
  first-order form \cref{eq:time-recast}, eliminate $\vv$ by differentiating the
  top equation and substituting, and recover the original second-order system
  \cref{eq:time-secondorder}. Conclude that the recast adds no physics---it only
  changes the shape of the unknowns.
  \item \emph{The explicit step-size bound.} For the scalar test $\dot y =
  -\lambda y$ with $\lambda = 50$, find the largest $\Delta t$ for which explicit
  Euler is stable. If a mesh refinement raises $\lambda_{\max}$ by a factor of $4$,
  by what factor must $\Delta t$ shrink? What does this say about explicit methods
  on fine meshes?
  \item \emph{Backward Euler is A-stable.} Show that for $\dot y = -\lambda y$ with
  \emph{any} complex $\lambda$ in the left half-plane ($\Re\lambda > 0$, decaying)
  and \emph{any} $\Delta t > 0$, the backward-Euler factor $g = 1/(1 +
  \lambda\Delta t)$ satisfies $|g| < 1$. (Hint: $|1 + \lambda\Delta t|^2 = (1 +
  \Delta t\,\Re\lambda)^2 + (\Delta t\,\Im\lambda)^2$.)
  \item \emph{Map versus fold, by the type.} Explain why \code{solve\_family}
  (electrostatics) can be evaluated on three machines in any order but
  \code{fold\_solve} (transient) cannot, citing the one feature of the
  \code{fold\_solve} signature that forbids reordering. Which line of the signature
  is it?
  \item \emph{Checkpoint exactness.} Using the schedule-split law
  \cref{eq:schedule-split}, argue that a transient run checkpointed after every
  $K$ steps and resumed produces \emph{exactly} the same final state as an
  uninterrupted run, for any $K$. What single quantity must the checkpoint save?
  Why does an electrostatic \code{solve\_family} have no analogous ``resume''?
  \item \emph{One implicit step by hand.} For the two-degree-of-freedom system of
  \cref{wex:twodof}, take \emph{one} backward-Euler step of size $\Delta t = 0.2$
  (rather than $0.1$) from rest. Form the step matrix $\mat A = \mat I + \Delta
  t\,\Cop + \Delta t^2\Kop$, write the right-hand side, and solve the $2\times 2$
  system for $\vv_1$. Confirm $\mat A$ depends only on $\Delta t$ and the operators,
  not on the step index---the operator-capture-once property.
  \item \emph{Numerical damping, qualitatively.} Generalized-$\alpha$ with
  $\rho_\infty = 0$ annihilates the highest-frequency modes in one step, while
  $\rho_\infty = 1$ preserves them. For a transient simulation polluted by ringing
  on the smallest mesh elements (unresolved high modes), which setting would you
  choose, and what is the risk of choosing $\rho_\infty$ too small (too much
  damping)?
  \item \emph{The two routes to S-parameters.} The driven analysis computes a
  frequency response by sweeping one frequency per solve (a \code{frequency\_sweep}
  map); the transient analysis computes the same response by one pulsed march
  followed by a Fourier transform. Give one reason to prefer each route. (Consider:
  how many frequencies you want, and whether the response is needed at a few sharp
  frequencies or across a broad band.)
\end{enumerate}
